\documentclass[notitlepage]
\usepackage{graphicx} % Required for inserting images
\usepackage{amssymb}
\usepackage[magyar]{babel}
\usepackage[T1]{fontenc}

\title{Algoritmikus Játékelmélet}
\author{Ádám Gutbrod, Dorottya Ács}
\date{2025}
\begin{document}




\chapter [Kombinatorikus játék, éles és betli játék]{Kombinatorikus játék fogalma, éles játék, betli játék, nyerő stratégia. Irányított gráf magja, mag egyértelműsége DAG-ban. Játékok, ill. pozíciók típusa}


{Kombinatorikus játék fogalma}

Def.: Kombinatorikus játék: $ J=(G,N_W,N_L,N_T,V_0) $, ahol $ G, DAG, N_W,L,T$ nyelők partíciója, $v_0 \in V(G)$ kiinduló állapot \\

{Éles játék}
Éles játék: $N_L \cup N_T=\varnothing$

Betli játék: $N_W \cup N_T = \varnothing$ \\
Megf.: Ha $N_T = \varnothing \Rightarrow \forall$ kombinatorikai játék v.vhető éles játékra.

Biz.: Bevezetünk egy extra csúcsot és $\forall N_L$-beli csúcsból 1-l ir. élt ide \\

Tétel: Tetszőleges éles kombinatorikus játékban $\forall$ csúcshoz egyértelműen eldönthető, hogy I vagy II típusú.
\begin{itemize}
    \item I-es típusú csúcsból indulva a kezdő garantálni tudja a győzelmet
    \item II-es típusú csúcsból indulva a másodiknak lépő játékos tudja garantálni azt
\end{itemize}

Biz.: Fordított topologikus sorrendben meghatározzuk $\forall$ csúcs típusát
$v_i típusa$ ...
\begin{itemize}
    \item II-es típusú, ha $\nexists v_iv_j \in E(G)$, amire $v_j$ típusa II-es
    \item  I-es, ha $\exists v_iv_j \in E(G)$, amire $v_j$ II-es
\end{itemize}

Megfigyelés: 

\end{document}